\section*{Part V: Time, Competition, and Pricing Experiments}

\subsection*{33. Dynamic Pricing Starts with Perishable or Changing Capacity}

An airline seat after departure, a hotel room after a night passes, and a delivery slot after its time window closes cannot be stored and sold tomorrow. When capacity is perishable and demand fluctuates, adjusting prices over time can improve utilization and contribution. The same logic can apply to appointments, events, logistics, electricity, and some cloud services. However, dynamic pricing is not inherently appropriate for every category.

The founder should identify what is changing: actual capacity scarcity, cost of supply, urgency, forecast demand, inventory risk, or the value of the next slot. Price movement must improve the allocation of a scarce resource without creating disproportionate harm or surprise. A studio with a fixed number of genuinely intensive review places might use transparent enrollment windows; pretending the same digital course is scarce every Tuesday is not an economic argument.

\subsection*{34. Revenue Management: Protect Capacity for Higher-Value Demand}

Suppose a 20-seat event has 12 seats sold early. A late buyer might value certainty and pay more, but unsold capacity becomes worthless at the start. The problem is choosing when to accept low-price demand versus holding capacity for uncertain later demand. Airlines and hotels use revenue-management systems because booking patterns, cancellation risk, and perishable capacity make that trade-off material.

In a tiny founder-led course, sophisticated yield management may be unnecessary and harmful to trust. A simpler approach is to set a transparent founding or early-commitment fee that compensates the buyer for earlier uncertainty, and a clearly disclosed later price with unchanged scope. Measure whether late demand actually exists rather than copying airline tactics because they sound mathematically advanced.

\subsection*{35. Surge Pricing and Fairness Constraints}

If demand spikes during a storm, a dynamic price can ration scarce supply and encourage additional capacity. It can also feel exploitative when people have no safe alternative. The economics of scarcity do not eliminate ethical and regulatory constraints. Fairness perceptions depend on the occasion, the seller's reason, local law, transparency, and the vulnerability of buyers.

Before adopting dynamic pricing, ask whether customers understand the rule, can anticipate the range, can decline without unacceptable harm, and whether pricing attracts additional supply or merely extracts more from those trapped. A long-lived brand should be able to explain its price changes in public.

\subsection*{36. Geographic, Seasonal, and Contractual Variation}

A company operating in multiple countries faces different purchasing power, taxes, currencies, distribution costs, competitors, and legal regimes. A single global price can create exclusion or invite arbitrage. Country-specific pricing may be rational, but must account for cost to serve, collection risk, foreign exchange, cross-border resale, and public perception.

Seasonal pricing is appropriate when demand and costs genuinely fluctuate. Contractual variation may reflect implementation scope, security requirements, service level, duration, or volume. The founder should distinguish \emph{economic reasons for variation} from arbitrary opportunism. A large contract that requires additional compliance and staff may be more costly and valuable than the same nominal service sold to a small business.

\subsection*{37. The Competitor's Price Is Evidence, Not an Instruction}

A rival cutting price by half might have superior costs, external subsidies, less service, a different buyer, a temporary promotion, or an unsustainable acquisition strategy. Matching the number without understanding the business is dangerous. You may need to improve your costs, redesign scope, demonstrate an actual difference, or concede that the rival now provides a better choice.

Pricing power survives when customers can articulate a renewed reason to choose you at the realized price. It need not mean always increasing prices. A low-cost producer can have strong economics and weak premium pricing. A specialized high-touch service can command high prices and low scale. Neither model is inherently superior; the choice should match mission, resources, and buyer jobs.

\subsection*{38. Pricing Changes Who Enters the Product}

A price increase may filter casual purchasers or exclude capable buyers who lack liquidity. It may also attract customers who expect concierge access, raise objection intensity, and increase review burden. A discount may increase attendance from price-constrained serious learners or attract poorly matched participants who overestimate their available time. The direction of these selection effects cannot be assumed.

Observe segment composition, entry skill, utilization, outcomes, complaints, support, and referrals by price and package. A short-run increase in revenue per customer may deteriorate long-run institutional quality. For education, readiness assessments and clear effort requirements may select for seriousness more fairly than a high fee alone.
